\newpage
\section{Introduction}
\label{sec:introduction}

\subsection{Contexte du projet}
\paragraph{} Le projet "Tu n'y peux rien" s'inscrit dans le cadre de l'UE Génie Logiciel de deuxième année à CY Cergy Paris Université. L'objectif est de développer un jeu de cartes multijoueur en Java, respectant des règles spécifiques et intégrant des bots de différents niveaux de difficulté.

\paragraph{} Ce projet s'inscrit dans la continuité des enseignements de programmation orientée objet et de génie logiciel. Il mobilise des compétences variées : conception UML, développement Java, gestion de projet avec Git, tests unitaires et documentation.

\subsection{État de l'art des jeux de cartes}
\paragraph{} Les jeux de cartes occupent une place importante dans l'histoire du jeu vidéo. Des titres comme \textit{Uno}, \textit{Skip-Bo} ou \textit{7 Wonders} proposent des mécaniques basées sur l'enchère et la pose de combinaisons. Cependant, peu de jeux intègrent une règle de "carte 2 écrase tout" ni un système de bombes et jokers comme dans notre projet.

\paragraph{} Le jeu "Tu n'y peux rien" se distingue par plusieurs originalités. Tout d'abord, la carte 2 possède un statut particulier : elle écrase toute carte simple posée précédemment. Ensuite, les Jokers sont polymorphes : ils peuvent remplacer n'importe quelle carte manquante dans une combinaison. Enfin, trois niveaux de difficulté pour les bots offrent une expérience de jeu adaptée à tous les publics.

\subsection{Objectif du projet}
\paragraph{} Recréer ce jeu de cartes au format numérique où un joueur humain affronte 2 à 4 bots. Le jeu doit gérer les combinaisons (simple, double, série, bombe, double joker), les cartes spéciales (2, Joker), et calculer les scores en fin de partie. L'interface graphique doit être intuitive et responsive, avec un design soigné (thème casino).

\subsection{Organisation du rapport}
\paragraph{} La section \ref{sec:specification} présente le cahier des charges et les fonctionnalités attendues. La section \ref{sec:conception} détaille la conception du logiciel et la section \ref{sec:oeuvre} présente sa mise en œuvre. La section \ref{sec:manuel} est un manuel utilisateur. La section \ref{sec:deroulement} décrit le déroulement du projet. La section \ref{sec:conclusion} conclut et propose des améliorations.
